%% Supplementary Note: grammar of the Graph Lens Lite query language.
%% Standalone document; build with:  latexmk -pdf dsl_grammar  (in this folder)
%%
%% AI-DRAFTED -- REVIEW
%% Derived from src/managers/query.js (encodeQuery / decodeQuery / QueryAST).
%% The four examples were checked by running them through QueryManager.encodeQuery,
%% decodeQuery and QueryAST (vitest, jsdom). Re-check if query.js changes.

\documentclass[10pt]{article}
\usepackage[T1]{fontenc}
\usepackage[a4paper,margin=2.5cm]{geometry}
\usepackage{listings}
\usepackage{url}

\lstset{basicstyle=\ttfamily\small, columns=fullflexible, keepspaces=true,
  breaklines=true, frame=single, framesep=4pt, upquote=true}

\title{Supplementary Note: Grammar of the Graph Lens Lite query language}
\date{}

\begin{document}
\maketitle

The query language selects or filters nodes and edges by conditions on their
properties. The grammar below is written in ISO EBNF and describes the queries
that the editor accepts as valid and that the evaluator assigns a meaning to.
It is derived from the implementation in \texttt{src/managers/query.js}, which
tokenizes the query text with regular expressions and decodes the tokens into
nested groups.

\section*{Grammar}

\begin{lstlisting}
query          = ws , expression , ws ;
expression     = clause
               | group , { ws , connector , ws , group } ;
group          = "(" , ws , expression , ws , ")" ;

clause         = property , " " , condition ;
property       = scope , "::" , name , "::" , name ;
scope          = "Node filters" | "Edge filters" ;
name           = name char , { name char } ;
name char      = ? any character except ":" ? ;

condition      = range | inverted range | membership | predicate ;
range          = "BETWEEN" , ws1 , number , ws1 , "AND" , ws1 , number ;
inverted range = "LOWER THAN" , ws1 , number , ws1 ,
                 "OR GREATER THAN" , ws1 , number ;
membership     = "IN" , ws , "[" , value , { "," , value } , "]" ;
value          = { value char | "\" , ? any character ? } ;
value char     = ? any character except "]", "\" and "," ? ;
predicate      = "IS TRUE" | "IS FALSE" | "IS MISSING" | "IS FOREIGN" ;

number         = [ "-" ] , digits , [ "." , digits ] ;
digits         = digit , { digit } ;
digit          = "0" | "1" | "2" | "3" | "4" | "5" | "6" | "7" | "8" | "9" ;

connector      = "AND" | "OR" | "NOT" ;      (* case-insensitive *)
ws             = { white space } ;
ws1            = white space , ws ;
white space    = ? space, tab, line break or no-break space ? ;
\end{lstlisting}

\section*{Lexical rules}

\begin{itemize}
  \item \textbf{Property path.} A property is addressed as
    \texttt{Section::Group::Property}. The section is \texttt{Node filters} or
    \texttt{Edge filters}; the group is the last bracketed label of the input
    column (\texttt{Uncategorized Properties} when none is given); the property
    is the rest of the column name. Group and property names may contain spaces but not
    colons. Exactly one white-space character separates the property path from
    its condition keyword.
  \item \textbf{Keywords.} Condition keywords (\texttt{BETWEEN},
    \texttt{AND} inside a range, \texttt{LOWER THAN}, \texttt{OR GREATER THAN},
    \texttt{IN}, \texttt{IS TRUE}, \texttt{IS FALSE}, \texttt{IS MISSING},
    \texttt{IS FOREIGN}) are upper case with single spaces as shown.
    Connectors (\texttt{AND}, \texttt{OR}, \texttt{NOT}) are matched
    case-insensitively.
  \item \textbf{Numbers.} Decimal with optional leading minus and optional
    fractional part (\texttt{-1.5}, \texttt{10}); exponents and a leading
    decimal point (\texttt{1e3}, \texttt{.5}) are not accepted.
  \item \textbf{Category values.} Values in \texttt{IN [\ldots]} are taken
    verbatim, including surrounding white space (\texttt{IN [a, b]} lists
    \texttt{a} and \texttt{\textvisiblespace b}). A backslash escapes the next
    character, so \texttt{[}, \texttt{]}, \texttt{,}, \texttt{(}, \texttt{)}
    and \texttt{\textbackslash} can appear in a value
    (\texttt{IN [a\textbackslash,b]} lists the single value \texttt{a,b}).
  \item \textbf{Connectors need parentheses.} A connector is recognized only
    between a closing and an opening parenthesis, \texttt{(\ldots)~AND~(\ldots)}. Two clauses joined without parentheses are rejected. A single
    clause may stand alone or in parentheses.
\end{itemize}

\section*{Evaluation}

\begin{itemize}
  \item \textbf{Order.} A chain $g_1\ c_1\ g_2\ c_2\ g_3 \ldots$ is evaluated
    strictly left to right, without operator precedence:
    \texttt{(A) OR (B) AND (C)} means \texttt{((A) OR (B)) AND (C)}.
    Parentheses nest to any depth and determine grouping.
  \item \textbf{Connectors.} \texttt{AND} and \texttt{OR} are the Boolean
    conjunction and disjunction. \texttt{NOT} is binary:
    \texttt{(A) NOT (B)} means \texttt{A AND NOT B}. There is no unary
    \texttt{NOT}; \texttt{NOT (A)} is rejected.
  \item \textbf{Range.} \texttt{BETWEEN a AND b} holds for numeric values $v$
    with $a \le v \le b$. \texttt{LOWER THAN a OR GREATER THAN b} holds for
    numeric values with $v \le a$ or $v \ge b$. Both bounds are inclusive.
  \item \textbf{Membership.} \texttt{IN [\ldots]} holds when the value equals
    one of the listed values exactly (case-sensitive). A value holding several
    entries separated by \texttt{|} matches if any entry does.
  \item \textbf{Boolean predicates.} \texttt{IS TRUE} matches
    \texttt{true}/\texttt{TRUE}/\texttt{1} and \texttt{IS FALSE} matches
    \texttt{false}/\texttt{FALSE}/\texttt{0}, as text or number.
  \item \textbf{Missing values and element type.} A node is tested only
    against \texttt{Node filters} properties and an edge only against
    \texttt{Edge filters} properties. A clause on the other element type's
    property, or on an absent value, is false, with two exceptions:
    \texttt{IS MISSING} is true when the value is absent or empty or the
    property belongs to the other element type, and \texttt{IS FOREIGN} is true
    only in the latter case. Both predicates are inserted by the filter panel's
    AND join (non-strict and ``complete cases'' mode, respectively) and are
    rarely written by hand.
  \item \textbf{No constraint.} When the filter panel has no active filter,
    the generated query is empty and every element matches; typed text in the
    editor must be non-empty to be applied.
\end{itemize}

\section*{Validation}

The editor re-validates on every change. A query is invalid, and cannot be
applied, when it is empty, contains an empty group \texttt{()}, two adjacent
groups without a connector \texttt{)(}, an unmatched parenthesis, or any text
not covered by the grammar. Property paths that are syntactically valid but
absent from the loaded graph are highlighted without blocking the query; they
match no element.

\section*{Examples}

The following queries use the columns of the case-study data set
(Supplementary datafile S1) and were checked against the implementation.

\begin{enumerate}
  \item Nodes strongly regulated in large cysts (inverted range, bounds
    inclusive):
\begin{lstlisting}
Node filters::Bulk TX data::Large cysts [log2FC] LOWER THAN -1 OR GREATER THAN 1
\end{lstlisting}
  \item Membrane or secreted proteins with a known direction of Wnt pathway
    modulation:
\begin{lstlisting}
(Node filters::Subcellular location::SubLoc coarse IN [Membrane,Secreted]) AND (Node filters::Mechanisms::Wnt signaling IN [modulates positive,modulates negative])
\end{lstlisting}
  \item Nodes in either drug interference signature, excluding those without a
    bulk expression value for large cysts (nesting and binary \texttt{NOT}):
\begin{lstlisting}
((Node filters::Drug interference::tolvaptan IS TRUE) OR (Node filters::Drug interference::metformin IS TRUE)) NOT (Node filters::Bulk TX data::Large cysts [log2FC] IS MISSING)
\end{lstlisting}
  \item A mixed edge and node query: physical protein--protein interactions
    together with nodes whose ADPKD literature evidence value is at least 100:
\begin{lstlisting}
(Edge filters::Physical interactions::Protein-protein interaction IS TRUE) OR (Node filters::Literature evidence::ADPKD literature evidence BETWEEN 100 AND 911)
\end{lstlisting}
\end{enumerate}

The group of a property is the last bracketed label of its input column, so the
column \texttt{Large cysts [log2FC] [Bulk TX data]} is addressed as
\texttt{Node filters::Bulk TX data::Large cysts [log2FC]}.

\end{document}
